% One-part-buffer replacement for the midpoint profile in Section 5.
% Inserted after the continuous root-gap calculation and before Section 6.

Choose the one-part-buffer integer
\begin{equation*}
  k_{\mathrm{co}}
  :=\left\lceil r_4^{\mathrm{co}}\right\rceil+1,
  \qquad
  d_n:=k_{\mathrm{co}}-r_4^{\mathrm{co}}.
  \tag{5.13}
\end{equation*}
Then
\[
  1\le d_n<2.
\]
The root separation in (5.12) tends to infinity, so
$k_{\mathrm{co}}<r_+$ for all sufficiently large $n$.

This placement retains the complete leading root separation.  The additional
class is not used to create a macroscopic first-moment buffer; it supplies the
smallest phase-uniform positive displacement from the real signed root.  That
constant displacement is enough because the slope at the root is of order
$(\log n)^2$.

We first verify that the complete interval
$[r_4^{\mathrm{co}},k_{\mathrm{co}}]$ remains in the common root corridor.  If
$k$ lies in this interval, then $k,r_4^{\mathrm{co}}=\Theta(n/\log n)$ and
\[
  \left|\frac nk-\frac{n}{r_4^{\mathrm{co}}}\right|
  =
  \frac{n|k-r_4^{\mathrm{co}}|}{k r_4^{\mathrm{co}}}
  =O\!\left(\frac{(\log n)^2}{n}\right)
  =o\!\left(\frac{\log\log n}{\log n}\right).
\]
Hence the normalized slope estimate (3.7) applies throughout the interval.
Writing
\[
  \Phi_n(k):=L_{S_4}(n,k)+(\log 2)k,
\]
there are phase-independent constants $0<c_{\mathrm s}<C_{\mathrm s}$ such
that, for all sufficiently large $n$,
\[
  c_{\mathrm s}(\log n)^2
  \le \Phi_n'(k)
  \le C_{\mathrm s}(\log n)^2
  \qquad
  (r_4^{\mathrm{co}}\le k\le k_{\mathrm{co}}).
\]
Since $\Phi_n(r_4^{\mathrm{co}})=0$ and $1\le d_n<2$, integration gives
\begin{equation*}
  c_{\mathrm s}(\log n)^2
  \le \Phi_n(k_{\mathrm{co}})
  \le 2C_{\mathrm s}(\log n)^2.
  \tag{5.13a}
\end{equation*}

At this exact $k_{\mathrm{co}}$, let $p_i^{(n)}$ be the finite-$n$ optimizer
of $L_{S_4}$.  Let $\lambda_{n,\mathrm{def}}$ be the Lagrange multiplier for
the deficit-mean constraint, with sign convention
$e^{\lambda_{n,\mathrm{def}}i}$.  The Lagrange equations and the decomposition
(3.11) give
\begin{equation*}
  p_i^{(n)}
  \propto d_{\alpha-i}^{-1}e^{\lambda_{n,\mathrm{def}}i}
  \propto e^{\widehat\lambda_n i+h_n(i)},
  \qquad
  \widehat\lambda_n:=B_n+\lambda_{n,\mathrm{def}}.
  \tag{5.14}
\end{equation*}
The target deficit mean is
$\alpha-n/k_{\mathrm{co}}$.  The root corridor and the preceding
$O((\log n)^2/n)$ displacement show that this target differs from $T_0$ by
$O(\log\log n/\log n)$.  The finite-dual package therefore gives uniform
convergence to the four-point limiting optimizer and, after decreasing one
constant if necessary,
\begin{equation*}
  \min_{2\le i\le5}p_i^{(n)}\ge c_p>0
  \tag{5.15}
\end{equation*}
throughout the complete phase.

Round $k_{\mathrm{co}}p_i^{(n)}$ to preliminary integers
$\widetilde k_i$.  Define
\[
  e_0:=\sum_i\widetilde k_i-k_{\mathrm{co}},
  \qquad
  e_1:=\sum_i i\widetilde k_i-(\alpha k_{\mathrm{co}}-n).
\]
Then $e_0,e_1=O(1)$.  Apply the bounded correction
\begin{equation*}
  \Delta k_2=e_1-3e_0,
  \qquad
  \Delta k_3=2e_0-e_1,
  \qquad
  \Delta k_4=\Delta k_5=0.
  \tag{5.16}
\end{equation*}
It satisfies
\[
  \sum_i\Delta k_i=-e_0,
  \qquad
  \sum_i i\Delta k_i=-e_1,
\]
so the corrected vector obeys the two exact conservation laws.  Its
coordinates remain positive by (5.15).

Let
\[
  \mathbf k^*:=(k_{\mathrm{co}}p_i^{(n)})_{i=2}^{5},
  \qquad
  \Delta^{\mathrm{tot}}:=\mathbf k-\mathbf k^*.
\]
The displacement has bounded norm and satisfies both homogeneous constraint
equations.  It is therefore tangent to the feasible affine plane.  At the
optimizer the linear term vanishes, and the entropy Hessian is
$-\operatorname{diag}(1/k_i)$.  Consequently
\[
  L_{\mathbf k^*+\Delta^{\mathrm{tot}}}
  -L_{\mathbf k^*}
  =O\!\left(
    \sum_i\frac{(\Delta_i^{\mathrm{tot}})^2}{k_i^*}
  \right)
  =O(1/k_{\mathrm{co}}).
\]
We have obtained the exact integer profile
\begin{equation*}
  \mathbf k=(k_2,k_3,k_4,k_5),
  \qquad
  u_i=\alpha-i,
  \qquad
  k_i=\Theta(n/\log n).
  \tag{5.17}
\end{equation*}

Its exact signed first moment is
\begin{equation*}
  \mathbb E Z_{\mathbf k}^{\mathrm{sgn}}
  =
  2^{k_{\mathrm{co}}}
  \frac{n!}{\prod_i(u_i!)^{k_i}k_i!}
  2^{-\sum_i k_i\binom{u_i}{2}}.
  \tag{5.18}
\end{equation*}
Stirling's formula, the tangent estimate above, and (5.13a) give one
phase-uniform logarithmic margin: there are constants
$0<c_{\mathrm{nr}}<C_{\mathrm{nr}}$ such that
\begin{equation*}
  c_{\mathrm{nr}}(\log n)^2
  \le
  \log\mathbb E Z_{\mathbf k}^{\mathrm{sgn}}
  \le
  C_{\mathrm{nr}}(\log n)^2
  \tag{5.19}
\end{equation*}
for all sufficiently large $n$.  Indeed, the exact integer-profile exponent
differs from $\Phi_n(k_{\mathrm{co}})$ by $O(\log n)$, which is absorbed by
the $(\log n)^2$ margin.

This is the only scale change relative to midpoint placement.  The
partial-diagonal central range uses only that
$k_{\mathrm{co}}^{-1}\log\mathbb E Z=o(1)$, while the full corner needs the
reciprocal first moment to dominate the polynomial number of residual
profiles.  Both follow from (5.19):
\[
  \frac{(\log n)^2}{k_{\mathrm{co}}}=o(1),
  \qquad
  (k_{\mathrm{co}}+1)^4
  e^{-c_{\mathrm{nr}}(\log n)^2}=o(1).
\]
Every $u_i$ tends to infinity, so the independent and complete declarations
of a class are disjoint.  Therefore
$Z_{\mathbf k}^{\mathrm{sgn}}>0$ implies a genuine cocoloring with exactly
$k_{\mathrm{co}}$ classes.

Finally, the chromatic threshold of Section~4 satisfies
$k_\chi^-=r_++O(\log n)$, whereas
$k_{\mathrm{co}}=r_4^{\mathrm{co}}+O(1)$.  Since
$\log n=o(n/(\log n)^3)$, equation (5.11) yields
\begin{equation*}
  k_\chi^- - k_{\mathrm{co}}
  =
  \left[
    \frac{(\log 2)^2}{4}A_4(\delta_n)+o(1)
  \right]
  \frac{n}{(\log n)^3}.
  \tag{5.20}
\end{equation*}
Thus the one-part buffer retains the full leading root separation.
